\section{Conditional Resident and Physical Evaluation}
\subsection{A Separate Optimized Software Cohort}
Resident completion includes state initialization, cost loads, all waves, canonical/directory retirement, output writes and acknowledgment drain. Host preparation and full readback are separate. This boundary differs from native E2 and Process. We reuse an affinity/NUMA-bound packed CPU implementation on dual Xeon Gold 6530 processors, testing 8, 32 and 64 workers. The 64-worker setting minimizes mean resident time on both full 40-sample workloads. The allocation has 64 physical cores and a 55-CPU-time quota; 40 samples permit at most 40 concurrent solves. This is the strongest measured CPU here.

Three independent processes each execute one warmup and three formal repetitions. Table~\ref{tab:resident} reports the mean and sample SD of the three process means, not nine independent trials. H100 resident includes initialization, launches/control, completion reads and synchronization; full output transfer and research comparison are outside. The BNN GPU receipt lacks an execution-time executable digest; its source/build provenance is a documented limitation.

\begin{table}[t]\centering
\caption{Complete-batch resident software time (seconds). CPU dispersion is over three independent process means. This cohort is separate from native eight-core E2/Process.}\label{tab:resident}
\setlength{\tabcolsep}{4pt}
\begin{tabular}{lrr}\toprule & BNN40 & LSTM40\\\midrule
CPU (64 workers) & $0.699624\pm0.026647$ & $0.725510\pm0.021287$\\
Retained H100 mean & 0.151800 & 0.159498\\
\bottomrule\end{tabular}
\end{table}

\subsection{Finite Model: Assumed Frequency and Memory}
We evaluate a \emph{hypothetical 200 MHz core} with declared 32-byte memory service. Baseline latency/issue interval is 80/5 ns; a conservative sensitivity scenario uses 160/10 ns. These correspond to 16/1 and 32/2 cycles at the assumed frequency. Reads and writes share the 32 physical ports; finite queues, credits, return identity and blocking remain active. Issue ceilings of 204.8/102.4 GB/s are not measured HBM throughput; conservative is not worst-case.

\begin{table}[t]\centering
\caption{Complete finite-model results at \emph{assumed 200 MHz}; neither routed-frequency performance nor board measurement.}\label{tab:model}
\begin{tabular}{llrr}\toprule Case & Latency/issue & Cycles & Model s\\\midrule
BNN40 & 80/5 ns & 102,557,024 & 0.512785\\
LSTM40 & 80/5 ns & 104,154,414 & 0.520772\\
BNN40 & 160/10 ns & 135,128,351 & 0.675642\\
LSTM40 & 160/10 ns & 137,198,248 & 0.685991\\
\bottomrule\end{tabular}
\end{table}

These deterministic predictions have no sampling error bars. Baseline expansion consumes 77.09/77.15\% of BNN/LSTM cycles and finalization 20.95/20.89\%; all other mutually exclusive phases account for the remainder. Candidate throughput is about 25.4\% of the 64-lane ideal: finite scheduling, transport and publication constrain parallelism. Overlapping stall counters are not additive.

The conservative BNN prediction exceeds the fastest measured CPU process mean, 0.668921 s, despite lying below its three-process mean. Thus the conditional result does not establish robust CPU superiority; GPU is faster in both scenarios. Frequency changes alter service discretization and scheduling: $T=C(f,\text{service},\text{structure})/f$, with absolute nanoseconds rounded up to cycles. Dividing old cycles by a new frequency is invalid when that binding changes.

Resident predictions exclude host preparation/readback and native assembly. Packed descriptors are shared by selected CPU/FPGA representations. Host layout conversion is not board DMA; unrelated native cohorts cannot supply an end-to-end FPGA speedup.
